\verseline{}{然而王后早已被沉重的情念所伤，}
\verseline{}{让伤口在血脉里滋长，被暗火慢慢吞噬。\footnote{拉丁文以saucia、vulnus写狄多受伤，以caecus ignis写暗中燃烧的爱欲。狄多彻夜不能安睡、反复想起埃涅阿斯，也承接阿波罗尼俄斯《阿尔戈英雄纪》中美狄亚爱上伊阿宋后的描写。其中vulnus alit（“滋养伤口”）紧贴卢克莱修《物性论》4.1068的ulcus enim vivescit et inveterascit alendo（“伤口因滋养而复苏、日久加深”）：卢克莱修用这一病理比喻告诫人切断爱情的“饲料”，维吉尔反而让狄多主动“滋养”伤口。到本卷末，开头的比喻性“伤口”会成为真实的致命伤。}}
\verseline{}{那人的英勇时时回到她心中，他一族的}
\verseline{}{高贵也反复浮现；他的面容和话语牢牢}
\verseline{5}{钉在胸中，这忧念不让她的肢体安然休息。}
\verseline{}{次日，福玻斯的明灯巡照大地，}
\verseline{}{奥罗拉从天穹移开湿润的暗影，\footnote{“福玻斯的明灯”指太阳，“奥罗拉”是黎明女神；umentem umbram 是清晨仍带水汽的暗影。}}
\verseline{}{她神志不宁，对心意相通的妹妹这样说道：}
\verseline{}{“妹妹安娜，怎样的梦使我悬心惊惧！}
\verseline{10}{这位来到我们宫中的新客人究竟是怎样的人，}
\verseline{}{他面容何等不凡，胸怀何等刚毅，武力何等强盛！}
\verseline{}{我确信他出自神族，这信念绝非虚妄。}
\verseline{}{恐惧显出卑弱的心性。唉，他曾被怎样的}
\verseline{}{命运抛掷！他唱述自己熬过了怎样的战争！}
\verseline{15}{倘若我心中不曾立下坚定不移的决意，}
\verseline{}{不愿再以婚姻的纽带同任何人结合——}
\verseline{}{自从最初的爱用死亡欺骗了我；}
\verseline{}{倘若我不曾厌弃婚床和婚礼的火炬，}
\verseline{}{也许唯独对这一次过失，我会屈服。\footnote{狄多将屈服于埃涅阿斯称为 culpa。此词可指过失、罪责或有失体面的行为；在狄多的话中，它首先关涉她对亡夫叙凯乌斯的誓约以及王后身份所要求的 pudor。婚礼火炬是罗马婚仪的标志之一。}}
\verseline{20}{安娜，我要坦白：自从我可怜的丈夫}
\verseline{}{叙凯乌斯死去，家神被兄弟的屠杀溅上鲜血，\footnote{第一卷叙述，狄多的兄弟皮格马利翁在家中祭坛前杀死叙凯乌斯。此处 Penates 既可指家神，也可转指家宅与炉边；“被屠杀溅上鲜血”因此同时写到家中的谋杀和祭坛的亵渎。}}
\verseline{}{只有这个人使我的心意转变，推动我}
\verseline{}{动摇的心：我认出了旧日火焰的痕迹。}
\verseline{}{可我宁愿最深的大地先在脚下裂开，}
\verseline{25}{宁愿全能的父神用雷霆把我击入幽影，}
\verseline{}{厄瑞玻斯苍白的幽影和无底的黑夜，}
\verseline{}{贞节啊，也不愿侵犯你，解开你的法则。}
\verseline{}{最先同我结合的那个人带走了我的爱；}
\verseline{}{让他拥有它，同他一起保存在坟墓里。”}
\verseline{30}{她说完，涌出的泪水浸透了胸襟。}
\verseline{}{安娜回答：“姐姐啊，对我而言，你比光明更亲爱，}
\verseline{}{难道你要独自在悲哀中耗尽全部青春，}
\verseline{}{既不知儿女的甜美，也不知维纳斯的馈赠？}
\verseline{}{你以为尘灰或埋葬的亡魂还会在意这些吗？\footnote{cinis 是火葬后留下的尘灰，Manes 是死者的亡魂或地下神灵。安娜并列死者的物质遗存与幽冥存在，借此否认叙凯乌斯仍会关切狄多是否再婚； “埋葬的”在语法上同亡魂结合，意义则笼罩整个表达。}}
\verseline{35}{就算从前没有一个求婚者能使你病苦的心屈服，}
\verseline{}{无论利比亚的，还是早先在推罗的；伊阿耳巴斯受你轻视，}
\verseline{}{其他首领也是如此，他们生在这片}
\verseline{}{富于战功的非洲土地。如今你还要抗拒合意的爱吗？}
\verseline{}{你难道不曾想到，自己定居在怎样的土地上？}
\verseline{40}{一边是盖图利亚人的城邑，一个战无不胜的种族；}
\verseline{}{无马衔的努米底亚人和不容宾客的叙耳提斯环绕着你；}
\verseline{}{一边是干渴荒芜的地域，还有四处}
\verseline{}{横行的巴耳卡人。何必再说推罗兴起的战争，}
\verseline{}{和你兄弟的威胁？\footnote{伊阿耳巴斯是北非君王，也是曾被狄多拒绝的求婚者。安娜列举盖图利亚人、努米底亚人、叙耳提斯沙洲、巴耳卡人以及来自推罗的威胁，以此劝狄多接受婚姻，巩固新城。第四卷第 44 行“和你兄弟的威胁”是《埃涅阿斯纪》现存约五十余处未足六音步的诗行之一，语意完整而格律未竟。}}
\verseline{45}{我深信，是众神显示吉兆、朱诺从旁相助，}
\verseline{}{伊利乌姆的船只才乘风走上这条航路。}
\verseline{}{姐姐，你将看见这座城何等壮大，怎样的王国}
\verseline{}{会因这样的结合而兴起！有特乌克罗斯人的武力相随，}
\verseline{}{迦太基的荣光将凭何等伟业高高升起！}
\verseline{50}{你只须祈求众神恩准，献祭得蒙悦纳，}
\verseline{}{尽心款待他们，再编织种种理由挽留，}
\verseline{}{趁冬季和多雨的俄里翁仍在海上肆虐，}
\verseline{}{船只受损，天空仍不宜航行。”\footnote{俄里翁既是神话猎人，也是与季候相联的星座；它的升落在古代航海传统中常被视为风雨和海况恶劣的征兆。安娜把受损船只和冬季天气当作挽留埃涅阿斯的现实理由。}}
\verseline{}{这番话以更炽烈的爱点燃她的心，}
\verseline{55}{给犹疑的意念带来希望，也松开了羞耻的约束。}
\verseline{}{她们先去神庙，在一座座祭坛前寻求和解；}
\verseline{}{依礼宰杀选出的绵羊，献给}
\verseline{}{立法的刻瑞斯、福玻斯和父神吕埃乌斯，}
\verseline{}{尤其献给最关心婚姻纽带的朱诺。\footnote{“吕埃乌斯”是巴克科斯的称号，含“解放者”之意。刻瑞斯被称为立法者，福玻斯具有神谕之权；朱诺尤其主管婚姻，因此这里特别向她献祭。}}
\verseline{60}{最美的狄多亲自用右手托住祭酒盘，}
\verseline{}{把酒浇在一头雪白母牛的双角之间，}
\verseline{}{又在神像前沿堆满祭品的祭坛缓步而行}
\verseline{}{一整天不断奉上祭品；她凝视牲畜剖开的}
\verseline{}{胸膛，从仍在搏动的内脏中求问征兆。\footnote{狄多将祭酒浇在白牛双角之间，再察看剖开的胸腔与内脏，是献祭与观兆程序。spirantia exta指仍带生命气息的内脏。65行vatum可修饰mentes，指先知无知的心，常解为先知未认识爱欲的力量；也有注家将它接ignarae，读为不懂先知及其预言的心，联系狄多或两姐妹。}}
\verseline{65}{无知的先知之心啊！誓愿对狂乱的人有何用，}
\verseline{}{神庙又有何用？柔和的火焰正在吞噬骨髓，}
\verseline{}{一个无声的伤口还活在胸膛之下。}
\verseline{}{不幸的狄多燃烧着，狂乱地走遍}
\verseline{}{全城，如一头中了飞箭的母鹿；}
\verseline{70}{牧人在克里特的林中远远射中了毫无防备的她，}
\verseline{}{一路追猎，却不知道已经把带翼的铁镞}
\verseline{}{留在她身上；她逃遍狄克忒的森林和山谷，}
\verseline{}{致命的箭杆却始终钉在肋侧。\footnote{母鹿明喻使爱欲的伤口成为具体箭伤。牧人不知自己留下致命箭镞，可与埃涅阿斯在这段关系中所起的作用对照；狄克忒是克里特的山名。}}
\verseline{}{有时她带着埃涅阿斯穿行城中}
\verseline{75}{向他展示西顿的财富与已经备妥的城市；}
\verseline{}{她开始说话，声音却停在半途。}
\verseline{}{有时，白日渐近黄昏，她又设下同样的宴席，}
\verseline{}{神志昏乱，再次要求聆听伊利乌姆的苦难，}
\verseline{}{又一次悬心听着，目光不离讲述者的唇边。}
\verseline{80}{后来众人散去，昏暗的月亮又一次}
\verseline{}{收敛光辉，西沉的群星劝人入睡，}
\verseline{}{她独自在空荡的宫中悲哀，伏在被遗下的}
\verseline{}{卧榻上；两人相隔，她仍听见、看见远方的他；}
\verseline{}{或因孩子身上父亲的形貌所迷，她把阿斯卡尼乌斯}
\verseline{85}{留在怀里，仿佛能借此欺骗那不可言说的爱。\footnote{阿斯卡尼乌斯的面貌像父亲，狄多因而把他留在怀中，试图借此骗过爱欲。infandus兼有不可言说、可怕之意；此处“不可言说”也可使人想到她的亡夫誓约与名誉。}}
\verseline{}{已经动工的塔楼不再升起，青年不再操练}
\verseline{}{兵器，也不再建造港口和战时}
\verseline{}{安全的堡垒；工程中断，巨大的城垛}
\verseline{}{和齐天的起重架都悬在那里。\footnote{第一卷中狄多以建城者和立法者的身份出现；此处塔楼、港口、堡垒和起重架一齐停住，她的私人欲望已经妨碍城邦建设。后来特洛伊舰队重新活动，正与迦太基工程的停滞相对。}}
\verseline{90}{朱庇特亲爱的妻子一看出她受这灾害所制，}
\verseline{}{名声也已不能阻挡狂乱，}
\verseline{}{萨图尔努斯之女便用这样的话向维纳斯进逼：}
\verseline{}{“好一场非凡的荣誉，好一份丰厚的战利品，}
\verseline{}{你和你的儿子确实名声伟大，永当传扬——}
\verseline{95}{只因一个女人被两位神明合谋欺骗！}
\verseline{}{我也清楚，你一直忌惮我们的城墙，}
\verseline{}{把巍峨迦太基的宫室视作威胁。}
\verseline{}{可这较量何时才止，何必相争至此？}
\verseline{}{我们何不缔结永久和平，成全约定的婚礼？}
\verseline{100}{你已经得到全心追求的一切：}
\verseline{}{狄多燃烧在爱中，狂乱已透入骨髓。}
\verseline{}{那么让我们以同等的神兆共同治理这个民族；}
\verseline{}{就让她侍奉一位佛里吉亚丈夫，}
\verseline{}{让推罗人作为嫁资交到你右手。”\footnote{朱诺所提议的“共同统治”含有明确的支配关系：狄多将侍奉佛里吉亚丈夫，推罗人则作为嫁资交给维纳斯。神祇谈判把个人结合直接换算为民族、城市和统治权的转让；朱诺意在使埃涅阿斯永久滞留非洲，从而偏离意大利的命运。}}
\verseline{105}{维纳斯看出她怀着虚假的心意开口，}
\verseline{}{想把意大利的王国引离，转向利比亚海岸，}
\verseline{}{便这样回答：“谁会疯狂到拒绝这种提议，}
\verseline{}{宁愿同你在战争中较量——}
\verseline{}{只要命运相助，使你所说的事成真？}
\verseline{110}{但我对命运并无把握，不能确定朱庇特是否愿意}
\verseline{}{让推罗人与特洛伊来客共有一座城，}
\verseline{}{是否准许两个民族混合，缔结盟约。}
\verseline{}{你是他的妻子，有权以祈求试探他的心意。}
\verseline{}{去吧，我会跟随。”天后朱诺便接过话来：}
\verseline{115}{“这件事由我承担。现在请留心，我用几句话告诉你，}
\verseline{}{眼前之事可以怎样完成。}
\verseline{}{埃涅阿斯和最不幸的狄多准备一同}
\verseline{}{进森林狩猎，等明日的提坦初升，}
\verseline{}{用光芒重新揭开大地。}
\verseline{120}{猎队奔忙、用围网圈住山谷时，我会从上方}
\verseline{}{倾下一片夹杂冰雹的乌云，}
\verseline{}{以雷霆震动整个天空。}
\verseline{}{随从将四处奔散，被浓黑的夜色遮住；}
\verseline{}{狄多和特洛伊人的首领会来到同一座}
\verseline{125}{洞穴。我会在场；只要你的意愿坚定，}
\verseline{}{我便以不渝的结合把她交给他，使她归属他；\footnote{Conte 2019 将第 126 行以方括号排印；该行重现第一卷第 73 行，该版校勘记记录 Peerlkamp 与 Ribbeck 的删疑。}}
\verseline{}{这就将是婚礼。”库忒瑞亚没有反对她的请求，}
\verseline{}{点头应允，并为识破这诡计而微笑。\footnote{repertis既可解“被识破的”，也可解“被想出的”：维纳斯或因识破诡计而笑，或对朱诺想出的办法发笑（参见 Ingo Gildenhard，\href{https://doi.org/10.11647/OBP.0023}{\textit{Virgil, Aeneid, 4.1–299}}，Open Book Publishers，2012，评4.127—128）。第105—106行已说她看出朱诺要将意大利的王国转向利比亚；她明知朱诺另有用意，仍点头配合。}}
\verseline{}{与此同时，升起的奥罗拉离开了大洋。}
\verseline{130}{晨光一现，选出的青年便走出城门；}
\verseline{}{疏眼的罗网、猎网和阔刃的猎矛一齐出动，}
\verseline{}{马叙利亚骑兵与一大群嗅觉敏锐的猎犬奔涌而出。}
\verseline{}{布匿人的首领在门前等候王后，她仍在寝宫迟延；}
\verseline{}{一匹披着紫红和黄金的骏马站在那里，}
\verseline{135}{烈性勃发，咬着冒出白沫的马衔。}
\verseline{}{她终于走出，由大群随从簇拥，}
\verseline{}{身披一件彩绣镶边的西顿斗篷。}
\verseline{}{她佩着金箭袋，头发以黄金绾起，}
\verseline{}{金扣又把紫红衣袍系在身上。}
\verseline{140}{佛里吉亚的同伴和欢喜的尤卢斯也向前走去。}
\verseline{}{埃涅阿斯本人比众人都更俊美，}
\verseline{}{走到她身旁，把两支队伍合在一起。}
\verseline{}{正如阿波罗离开冬日的吕喀亚与克珊托斯河，}
\verseline{}{去探望母亲的提洛岛，}
\verseline{145}{重新举行舞蹈；克里特人、德律俄珀斯人与彩绘身体的}
\verseline{}{阿伽图尔西人混在一起，围着祭坛欢呼；}
\verseline{}{他亲自走在库恩图斯山脊上，用柔软的叶枝压住}
\verseline{}{流动的长发，把它理好，再用黄金缠束；}
\verseline{}{箭矢在肩头鸣响。埃涅阿斯走来并不稍逊于他，}
\verseline{150}{如此光彩从他出众的面容上闪耀。\footnote{埃涅阿斯被比作阿波罗离开吕喀亚与克珊托斯，前往母亲所在的提洛，重新举行祭舞；库恩图斯是提洛的山。这些地名都与阿波罗传说相连。各族人群在祭坛四周混杂欢呼，与狩猎队伍相映。}}
\verseline{}{他们抵达高山与无路可通的兽林以后，}
\verseline{}{看啊，野山羊从岩峰上被赶落，}
\verseline{}{沿山脊飞奔；另一边，鹿群疾驰着}
\verseline{}{穿越开阔的平野，在尘土飞扬的逃奔中聚成队列，}
\verseline{155}{把高山抛在身后。}
\verseline{}{少年阿斯卡尼乌斯骑着烈马，在山谷中央}
\verseline{}{满心欢喜，时而飞越这一群，时而超过那一群；}
\verseline{}{他祈望在这些迟钝的牲畜中，能有一头口吐白沫的}
\verseline{}{野猪送到眼前，或一头黄狮从山上下来。}
\verseline{160}{这时天空开始混乱，响起巨大的轰鸣；}
\verseline{}{紧接着乌云涌来，其中夹着冰雹。}
\verseline{}{推罗的随从、特洛伊的青年，}
\verseline{}{还有那位来自达耳达尼亚的维纳斯之孙，也因恐惧四散田野，}
\verseline{}{各自寻找遮蔽；山洪奔泻而下。}
\verseline{165}{狄多和特洛伊人的首领来到同一座}
\verseline{}{洞穴；首先，大地与主持婚礼的朱诺一同发出信号；}
\verseline{}{火光闪耀，知情的天空见证这场结合，}
\verseline{}{宁芙们在最高的山巅发出哀号。\footnote{洞穴结合出现互相冲突的征兆：大地和朱诺如同为婚礼作证，闪电或火光照亮天空，宁芙却在山巅“哀号”。ululare 既可用于仪式性呼喊，也带有悲鸣意味；这场洞穴结合从一开始便兼有婚礼仪式和死亡凶兆。}}
\verseline{}{那一天是死亡、也是灾祸最初的}
\verseline{170}{起因；狄多不再为体面或名声动摇，}
\verseline{}{也不再筹划一场隐秘的爱；}
\verseline{}{她称它为婚姻，用这个名字遮掩自己的过失。\footnote{第169—170行把这一天称为死亡与灾祸的起因，将洞穴结合连到狄多的结局。狄多以“婚姻”称这次结合，叙述者则称之为culpa（参19行）；339行埃涅阿斯否认曾举起婚礼火炬、进入这样的盟约，双方对结合的名分说法不同（参见 \href{https://doi.org/10.11647/OBP.0023}{Gildenhard 2012，评注4.169—172}）。}}
\verseline{}{流言立刻穿行于利比亚的各座大城——}
\verseline{}{流言，世上没有任何祸害比她更为迅疾；}
\verseline{175}{她因奔走而壮大，在行进中取得力量，}
\verseline{}{初起时身形微小，很快便升入高空；\footnote{Conte 2019 在第 176 行采用 Baehrens 的 initu（“在开始时”），主要传本作 metu（“在恐惧中”）；按正文读法，initu primo 即“初起时”。}}
\verseline{}{她踏在地上，头却藏入云层。}
\verseline{}{大地母亲因对众神的愤怒而激荡，}
\verseline{}{据说在科俄斯与恩刻拉多斯之后，最后又生下}
\verseline{180}{这个妹妹；她双足飞快，有敏捷的翅膀，}
\verseline{}{是个骇人而庞大的怪物。她身上有多少羽毛，}
\verseline{}{羽毛下面便有多少不眠的眼睛——说来也令人惊异——}
\verseline{}{多少舌头，多少张喧响的嘴，也竖起多少只耳朵。}
\verseline{}{夜间，她在天地之间的阴影里飞翔，}
\verseline{185}{发出尖厉的叫声，双眼从不屈服于甜美的睡眠；}
\verseline{}{白昼，她守望在最高的屋脊，}
\verseline{}{或高耸的塔楼上，使一座座大城惊恐；}
\verseline{}{她紧抓虚假与邪曲，一如传报真实。\footnote{Fama 同时是名声、传闻、舆论以及人格化的“流言”。维吉尔把她写成大地继科俄斯和恩刻拉多斯之后生下的巨怪；遍布羽毛、眼睛、舌头、嘴和耳朵的身体，使信息传播本身成为一种不断增殖的生物力量。她同样执着于真实与虚假，含义比“谣言”更为宽广。}}
\verseline{}{那时她欣喜地用纷繁的言语填满众人的耳朵，}
\verseline{190}{把发生过的与从未发生的事一同传唱：}
\verseline{}{特洛伊血统的埃涅阿斯已经到来，}
\verseline{}{美丽的狄多竟肯屈尊同这男人结合；}
\verseline{}{如今整个漫长的冬季，他们彼此沉溺于奢靡，}
\verseline{}{忘记自己的王国，被可耻的欲望攫住。\footnote{流言以事实为起点：两人确实共同度过时日，狄多的城市也确实停工；但“整个漫长的冬季”“沉溺奢靡”和“忘记王国”已经把事实组织成道德指控。Fama 的力量正在于把发生过的事与夸张、评价和未发生之事混合。}}
\verseline{195}{这丑恶的女神便把这些流言传遍众人之口。}
\verseline{}{旋即掉转行程，奔向伊阿耳巴斯国王，}
\verseline{}{用言语点燃他的心，又将愤怒层层堆高。}
\verseline{}{他是阿蒙从被掳的加拉曼提斯宁芙那里生下的儿子，\footnote{伊阿耳巴斯被说成朱庇特·阿蒙与一位加拉曼提斯宁芙之子。阿蒙是北非神祇，希腊—罗马传统常将其与朱庇特等同；这一谱系使伊阿耳巴斯对朱庇特的祷告兼有血缘诉求与地方宗教权威。}}
\verseline{}{在自己广阔的国土上为朱庇特建立了}
\verseline{200}{一百座宏伟神庙、一百座祭坛，又献出不眠的火，}
\verseline{}{作为众神永恒的守卫；祭牲的鲜血}
\verseline{}{使土地肥沃，各色花环盛放在门槛上。}
\verseline{}{他心神狂乱，被苦涩的传闻点燃，}
\verseline{}{据说站在祭坛前、众神的神像之间，}
\verseline{205}{仰起双手，以求告者的身份向朱庇特反复祈求：}
\verseline{}{“全能的朱庇特，如今毛鲁西亚人躺在彩绘的宴榻上，}
\verseline{}{向你奠下勒奈乌斯的礼酒，\footnote{“勒奈乌斯”是巴克科斯的称号，常与葡萄压榨和酒相联；此处“勒奈乌斯的礼酒”即宴饮中向朱庇特奠献的葡萄酒。伊阿耳巴斯以北非祭祀者的身份质问朱庇特是否仍观看并回应人间。}}
\verseline{}{你看见这些了吗？父神啊，难道你挥动雷霆时，}
\verseline{}{我们只是徒然颤栗？云中的暗火}
\verseline{210}{只是惊吓人心，发出空洞的轰鸣？}
\verseline{}{一个女人流浪到我们的疆界，出资建起}
\verseline{}{一座小城；我们给她海岸耕作，}
\verseline{}{也给她当地的法权，她却拒绝同我结合，\footnote{伊阿耳巴斯把狄多说成以钱取得土地并建立小城，暗合柏尔萨建城传说：狄多购得一张牛皮所能围住的土地，再把牛皮裁成细条扩大范围。祷词故意缩小她的建城功业，并把土地和法权说成当地男性君王的恩赐。}}
\verseline{}{把埃涅阿斯当作主人迎进自己的王国。}
\verseline{215}{如今那个帕里斯带着一群不男不女的随从，}
\verseline{}{用迈俄尼亚软帽的系带束住下巴与润香的头发，\footnote{伊阿耳巴斯把埃涅阿斯辱称为“帕里斯”，将他留恋外国王后同帕里斯拐走海伦相提并论。原文 semiviro comitatu 是“半男人的随从”，Maeonia mentum mitra 则以吕底亚式软帽、香发和系带制造东方奢靡与女性化的侮辱；这些都是伊阿耳巴斯用来羞辱埃涅阿斯的细节。}}
\verseline{}{享用他抢来的猎物；而我们还把礼物}
\verseline{}{带到你的神庙，守护一场空洞的名声。”}
\verseline{}{他这样祈祷，紧握祭坛，}
\verseline{220}{全能的神听见了，把目光转向王城的墙垣}
\verseline{}{和那两个忘却更高名声的爱人。}
\verseline{}{接着他这样对墨丘利说，传下命令：}
\verseline{}{“去吧，儿子，召来西风，乘着羽翼滑落；}
\verseline{}{达耳达诺斯人的首领如今在推罗人的迦太基}
\verseline{225}{徘徊，不再望向命运赐给他的城市——}
\verseline{}{你去告诫他，迅速穿过天空，带去我的话。}
\verseline{}{他最美的母亲曾许诺，他会成为另一番人物；}
\verseline{}{她两次从希腊刀兵下救他，原有更高目的；}
\verseline{}{她许诺他将统治怀着霸业、因战争沸腾的}
\verseline{230}{意大利，使一个民族从特乌克罗斯的高贵血脉中诞生，}
\verseline{}{让整个世界都服从它的法度。}
\verseline{}{如果如此伟业的荣光不能点燃他，}
\verseline{}{他自己也不肯为自己的荣耀承担劳苦，}
\verseline{}{身为父亲，他竟连罗马的城堡也不肯留给阿斯卡尼乌斯吗？}
\verseline{235}{他在筹划什么？抱着什么希望滞留于敌对的民族，}
\verseline{}{不再顾望奥索尼亚的子孙和拉维尼乌姆的田野？}
\verseline{}{让他航行。这就是全部。你就把这句话带去。”}
\verseline{}{他说完。墨丘利准备服从伟大父神的}
\verseline{}{命令；首先把金色的飞鞋系在脚上，}
\verseline{240}{那双鞋以羽翼托起他的身体，无论越过海面}
\verseline{}{还是陆地，都同迅疾的风一起带他前行。}
\verseline{}{接着他拿起神杖：凭它从奥尔库斯唤出}
\verseline{}{苍白的魂灵，把另一些送入幽暗的塔耳塔洛斯；}
\verseline{}{他赐下或夺去睡眠，也在死亡中启封人的眼睛。\footnote{朱庇特遣墨丘利催促埃涅阿斯离开迦太基的场景承接《奥德赛》5.28—42：宙斯遣赫尔墨斯前往卡吕普索处，命她放奥德修斯返航。墨丘利的神杖后来常称“商神杖”或caduceus，兼具引魂与司梦之权：从冥界召魂、送魂入地下、赐予或夺去睡眠。末句\textit{lumina morte resignat}可解作“在死亡中启封双眼”，联系火葬前重开死者双眼的习俗（见\href{https://penelope.uchicago.edu/Thayer/L/Roman/Texts/Pliny_the_Elder/11*.html}{老普林尼《博物志》11.150}）；塞尔维乌斯（4.244）则把resignat释为“合上”。}}
\verseline{245}{凭着神杖，他驱策长风，游过混乱的}
\verseline{}{云层；飞行间，他已经望见强悍的阿特拉斯}
\verseline{}{高峰与陡峭的山腹，那巨神用头顶支撑天空。}
\verseline{}{阿特拉斯长满松树的头颅不断被暗云}
\verseline{}{围住，经受风雨击打；}
\verseline{250}{倾落的雪覆盖双肩，河流从他的下巴}
\verseline{}{倒泻而下，老人蓬乱的胡须冻成坚冰。}
\verseline{}{库勒涅之神先在这里凭一双平衡的羽翼}
\verseline{}{停住，继而全身俯冲，直落向海浪，}
\verseline{}{像一只鸟贴近水面，绕着海岸、绕着}
\verseline{255}{多鱼的礁石低飞。}
\verseline{}{正是这样，他飞行在天地之间，}
\verseline{}{掠过利比亚的沙岸，切开长风；}
\verseline{}{库勒涅之神正从外祖父那里飞来。\footnote{阿特拉斯的身体被写成山岳：头颅负天，肩披积雪，河水从下巴流下，胡须冻结。墨丘利之母迈亚是阿特拉斯之女，因此本行称他“从外祖父那里”飞来； “库勒涅之神”则以阿卡狄亚的库勒涅山称呼墨丘利，那里是他的出生地和崇拜地。}}
\verseline{}{飞鞋一触及那些茅屋，}
\verseline{260}{他便看见埃涅阿斯正在奠立堡垒，更新房舍；}
\verseline{}{那人佩着宝剑，黄褐碧玉如星镶嵌其上，}
\verseline{}{从肩上垂落的斗篷燃烧着推罗紫；}
\verseline{}{富有的狄多制成这件礼物送给他，}
\verseline{}{又把纤细金线织进衣料。\footnote{埃涅阿斯所佩宝剑镶有黄褐碧玉（jasper），肩上斗篷染着推罗紫；这身装束已经带上狄多宫廷的奢华色彩。斗篷由狄多赠送并以金线织成；墨丘利恰在他参与修建迦太基时看见这一形象，随即以朱庇特的命令催他离开。}}
\verseline{265}{他立即进逼：“你如今竟在为巍峨的迦太基}
\verseline{}{奠立根基，依顺一个女人，建造她美丽的城，}
\verseline{}{唉，竟忘了自己的王国与本业？}
\verseline{}{是众神的统治者亲自从光辉的奥林匹斯}
\verseline{}{把我派来——他以神威转动天地；}
\verseline{270}{正是他命我飞过疾风，带来这道命令：}
\verseline{}{你在筹划什么？抱着什么希望消磨时日在利比亚？}
\verseline{}{如果如此伟业的荣光不能触动你，\footnote{Conte 2019 正文由第 272 行直接接第 274 行；该版校勘记注明传统第 273 行在多种重要传本中缺失。}}
\verseline{}{也该看看长大的阿斯卡尼乌斯，看看继承人尤卢斯的}
\verseline{275}{希望；意大利的王国和罗马的大地}
\verseline{}{都应归他。”库勒涅之神说完这些，}
\verseline{}{还在说话时，便离开凡人的眼前，}
\verseline{}{远远从眼前消失，化入稀薄的空气。}
\verseline{}{埃涅阿斯却被这景象震住，默然失神，}
\verseline{280}{他惊呆了，头发竖起，声音哽在喉中。}
\verseline{}{他迫切要逃走，离开这甜美的土地，}
\verseline{}{为神明如此重大的警告和命令所震骇。}
\verseline{}{唉，他该怎么办？如今敢用怎样的话走近}
\verseline{}{狂怒的王后？又该从哪里开口？}
\verseline{285}{他的心意迅疾，一时分向这里，一时分向那里，}
\verseline{}{又被扯往不同方向，把一切反复思量。}
\verseline{}{犹疑之间，这个决定看来更好：}
\verseline{}{他召来门涅斯透斯、塞耳革斯图斯和勇敢的塞雷斯图斯，}
\verseline{}{命他们暗中整备舰队，把同伴召集到岸边，}
\verseline{290}{备妥船具，却不说明为何突然改变计划；}
\verseline{}{而他自己，既然最善良的狄多毫不知情，}
\verseline{}{也想不到如此深厚的爱会被割断，}
\verseline{}{便打算寻找接近她的途径、最温和的}
\verseline{}{开口时机，以及怎样行事最为妥帖。众人立即}
\verseline{295}{欣然服从命令，执行他的吩咐。}
\verseline{}{可王后预先觉察了欺瞒——谁能欺骗一个爱着的人？}
\verseline{}{她最先捕捉到将来的行动，}
\verseline{}{即使一切安稳，也对一切怀着恐惧。同一个邪恶的流言}
\verseline{}{又给狂怒的她带来消息：舰队正在整装，准备启航。}
\verseline{300}{她心无所主地狂怒，燃烧着在全城}
\verseline{}{作酒神狂舞，如被摇动的圣物激起的提伊阿斯；\footnote{狄多被比作参加巴克科斯狂仪的提伊阿斯。隔年举行的祭典按包含首尾计算称“三年一度”；喀泰戎山既是底比斯酒神崇拜的重要地点，也是彭透斯悲剧的发生地。这个明喻把她的奔走写成受神力驱动的仪式性狂乱。}}
\verseline{}{一听见巴克科斯，三年一度的狂欢祭便驱策她，}
\verseline{}{夜间的喀泰戎山也用呼喊召唤她。}
\verseline{}{她终于主动用这些话向埃涅阿斯发难：}
\verseline{305}{“背信的人，你竟还指望能隐瞒}
\verseline{}{如此罪恶之事，悄无声息地离开我的土地？}
\verseline{}{我们的爱留不住你？从前交出的右手留不住你？}
\verseline{}{将以残酷方式死去的狄多也留不住你？}
\verseline{}{你甚至在冬日的星辰下整备舰队，}
\verseline{310}{急着在北风正盛时奔向海上，}
\verseline{}{残忍的人！假使你要去的并非异乡土地和}
\verseline{}{陌生家园，古老的特洛伊也依然存在，}
\verseline{}{难道你还会率领舰队越过波涛汹涌的大海去往特洛伊？}
\verseline{}{你竟是在逃离我吗？我凭这些眼泪和你的右手求你——}
\verseline{315}{因为可怜的我已没有给自己留下别的东西——}
\verseline{}{凭我们的婚姻，凭已经开始的婚礼，\footnote{conubia与hymenaei都属于婚姻、婚礼的语言；inceptos是“已开始的”。狄多以此强调双方结合的约束。她的弃妇控诉可与卡图卢斯第64首中阿里阿德涅责问忒修斯的段落参照；其中也并列conubia与hymenaeos。}}
\verseline{}{如果我曾对你有任何恩惠，我的任何事物}
\verseline{}{曾使你感到甜美，请怜悯这个倾颓的家——}
\verseline{}{我求你，若祈祷尚有余地——抛开这个念头。}
\verseline{320}{因为你，利比亚各族和努米底亚人的君王}
\verseline{}{憎恨我，推罗人也与我为敌；因为你，同样}
\verseline{}{熄灭的是我的贞节和从前的名声——唯有它曾把我送近星辰。}
\verseline{}{我将死去，你把我遗弃给谁？如今我只能叫你‘客人’，}
\verseline{}{曾经的丈夫只剩下这个称呼。\footnote{狄多先称埃涅阿斯为 hospes，即“客人”或受宾主法保护的异乡人，随即说“丈夫留给我的如今只有这个称呼”。狄多先把埃涅阿斯当作 coniux（“丈夫”），此刻又退回 hospes（“客人／宾主关系中的外来者”）；她以这一称呼责问他是否还守宾主信义。}}
\verseline{325}{我还等什么？等我的兄弟皮格马利翁摧毁城墙，}
\verseline{}{或等盖图利亚的伊阿耳巴斯把我俘去？}
\verseline{}{倘若在你逃走以前，我至少从你那里}
\verseline{}{怀上一个孩子，倘若宫中有一个小埃涅阿斯}
\verseline{}{玩耍，面容终究还能把你带回，}
\verseline{330}{我便不会觉得自己已被彻底俘获、彻底遗弃。”}
\verseline{}{她说完了。因朱庇特的警告，他始终保持}
\verseline{}{目光不动，勉力把痛苦压在心底。}
\verseline{}{终于，他简短回答：“王后，你所能历数的}
\verseline{}{种种恩德，我永远不会否认；}
\verseline{335}{只要我还记得自己，只要气息还支配这些肢体，}
\verseline{}{我也绝不会厌于记住厄利莎。}
\verseline{}{论到此事，我只说几句。我从未希望暗中}
\verseline{}{掩藏这次逃离——不要这样设想——也从未}
\verseline{}{举起丈夫的婚礼火炬，或走进这样的盟约。\footnote{婚礼火炬是婚仪的标志；foedus可指盟约、契约或婚姻结合。狄多称双方关系为婚姻，埃涅阿斯在此否认曾举起婚炬、进入这样的盟约。}}
\verseline{340}{如果命运允许我凭自己的权柄度过一生，}
\verseline{}{依照本愿安顿我所忧心之事，}
\verseline{}{我首先会守护特洛伊城，也珍重我亲爱的}
\verseline{}{族人遗民；普里阿摩斯高大的宫室仍会屹立，\footnote{Conte 2019 仅将 Priami tecta alta manerent（“普里阿摩斯高大的宫室仍会屹立”）置于方括号；该版校勘记记录 Peerlkamp 与 Ribbeck 的删疑。}}
\verseline{}{我也会亲手为战败者重建帕加马。}
\verseline{345}{但如今格律涅乌斯的阿波罗、吕喀亚的神谕}
\verseline{}{都命我追求伟大的意大利，意大利：}
\verseline{}{那里是我的爱，那里是我的祖国。如果你这腓尼基女人}
\verseline{}{能留恋迦太基的堡垒、利比亚城市的景象，}
\verseline{}{又为何不肯让特乌克罗斯人定居奥索尼亚的土地？}
\verseline{350}{我们也有权寻找异邦的王国。}
\verseline{}{每当夜色用潮湿的阴影覆盖大地，}
\verseline{}{每当燃烧的星辰升起，我父亲安喀塞斯}
\verseline{}{不安的形影便在梦中警告我，使我恐惧；}
\verseline{}{少年阿斯卡尼乌斯，还有我对这亲爱孩子的亏欠，也在告诫我：}
\verseline{355}{我正夺去他赫斯珀里亚的王国和命运赐予的田野。}
\verseline{}{如今连神明的传令者也由朱庇特亲自派来——}
\verseline{}{我凭你我二人的生命作证——穿过疾风带来命令；}
\verseline{}{我亲眼看见神明在显明的光中}
\verseline{}{进入城墙，亲耳听见他的声音。}
\verseline{360}{别再用你的怨诉把你我一同点燃；}
\verseline{}{我追寻意大利，并非出于本愿。”\footnote{意大利是“我的爱”和“我的祖国”，承接第三卷关于达耳达诺斯意大利祖源的说明；前往意大利，因而也是返回祖地。Italiam non sponte sequor说“我并不自愿追寻意大利”，把这次离去归于神命。}}
\verseline{}{他还在说话，她早已侧过脸去注视着他，}
\verseline{}{目光在这里那里翻转，又用沉默的眼睛}
\verseline{}{扫遍他全身，终于燃烧着开口：}
\verseline{365}{“女神怎会是你的母亲，达耳达诺斯怎会是你的族祖？}
\verseline{}{背信的人！可怖的高加索从坚硬岩石中生下你，}
\verseline{}{许耳卡尼亚的母虎才曾以乳汁哺育你。\footnote{高加索岩石和许耳卡尼亚母虎属于古代斥责冷酷者的传统意象：冷酷之人仿佛从荒蛮岩石诞生、由猛兽哺育。狄多借这套语言取消维纳斯与达耳达诺斯的神圣谱系，反转埃涅阿斯刚才诉诸的家世与命运。马克罗比乌斯（6.1.59）又把第 371—372 行狄多抱怨朱诺和朱庇特不再公正看待此事，与阿克齐乌斯失传悲剧《安提戈涅》中同样抱怨诸神不再治理、众神之王不再关怀的句子并列。}}
\verseline{}{可我何必再遮掩？还要留着自己等待什么更大的事？}
\verseline{}{他曾为我的哭泣叹息吗？曾转动目光吗？}
\verseline{370}{曾因屈服而流泪，或怜悯爱着他的我吗？}
\verseline{}{我该把哪一件放在前面？如今，正是如今，至高的朱诺}
\verseline{}{和萨图尔努斯之子也不再公正地看待这些事。}
\verseline{}{信义已无处安全。我接纳他——一个被抛上岸、}
\verseline{}{一无所有的人，疯狂地让他分享我的王国；}
\verseline{375}{我替他寻回失去的舰队，从死亡中救回同伴。}
\verseline{}{唉，我被复仇女神点燃，正被卷走！如今先是神谕者阿波罗，}
\verseline{}{又是吕喀亚的神谕，如今还有朱庇特亲自派来的}
\verseline{}{神使，穿过疾风带来骇人的命令。}
\verseline{}{这自然是天神的劳苦，是让安宁者不安的忧虑！\footnote{“这自然是天神的劳苦，是让安宁者不安的忧虑”具有强烈反讽。狄多假装承认天神必须为埃涅阿斯的航程劳碌，实际是在嘲讽他把离弃归因于天命和诸神的命令。}}
\verseline{380}{我不挽留你，也不反驳你的话。}
\verseline{}{去吧，乘风追随意大利，越过波浪寻找你的王国。}
\verseline{}{但我深信，如果公正的神明还有力量，你会在海中}
\verseline{}{暗礁上喝尽惩罚，一次次呼喊}
\verseline{}{狄多的名字。我虽不在场，也会以黑色火焰追随；}
\verseline{385}{等冰冷的死亡把灵魂从肢体上分开，}
\verseline{}{我的阴影会无处不在。恶人，你会付出代价。}
\verseline{}{我会听见；这消息也会传到地下最深处的我。”}
\verseline{}{话到这里，她猛然截断言辞，带着病恹恹的身躯}
\verseline{}{逃离天光，转身从他眼前离去，}
\verseline{390}{留下他因恐惧迟疑，准备说许多}
\verseline{}{话。侍女们扶住她，把瘫软的肢体}
\verseline{}{带回大理石寝宫，放在卧榻上。}
\verseline{}{虔敬的埃涅阿斯虽然想以安慰减轻她的痛苦，}
\verseline{}{用言语排解她的忧虑，}
\verseline{395}{虽不断叹息，深厚的爱也使他的心动摇，}
\verseline{}{仍然执行众神的命令，回到舰队。\footnote{叙述者在狄多最尖锐的控诉之后仍称 pius Aeneas。“虔敬”在此涵盖对神命、父亲、儿子以及特洛伊人未来的尽责。这个固定称号紧接狄多的痛苦出现：埃涅阿斯仍在履行使命，却也正在伤害狄多。}}
\verseline{}{这时特乌克罗斯人全都俯身劳作，从整个海岸}
\verseline{}{把高大的船只推下；涂油的龙骨浮在水上。}
\verseline{}{他们急于逃离，从林中运来还带枝叶的船桨}
\verseline{400}{和未经削制的原木。\footnote{特洛伊人为赶紧启航，从林中运来仍带枝叶的木材和未经修整的原木。诗中 frondentis remos 指仓促砍下、尚带叶枝而准备用作桨的木料；树枝尚未削净便拿来作桨，写出撤离之仓促。}}
\verseline{}{你会看见他们迁徙，从整座城里奔涌而出。}
\verseline{}{仿佛蚂蚁劫掠一大堆谷粒，}
\verseline{}{记得冬日将临，把它们贮进巢穴；}
\verseline{}{黑色的队伍穿过田野，沿一条狭窄小径}
\verseline{405}{把战利品运过草间；一部分用肩膀奋力}
\verseline{}{推动巨大的谷粒，一部分约束行列，}
\verseline{}{催促迟缓者；整条路都因劳作沸腾。\footnote{蚁群明喻沿用史诗中以微小生物显示集体规模的传统。谷粒被称作“战利品”，队伍中有搬运者、整队者和督促迟缓者；整个队伍像一支有序的军队，也像分工明确的城邦劳动者，这番忙碌又反衬狄多的失控。}}
\verseline{}{狄多啊，看见这些时，你怀着怎样的心情？}
\verseline{}{当你从城堡最高处远望广阔的海岸}
\verseline{410}{一片沸腾，看见整个海面在眼前}
\verseline{}{因巨大的呼喊混乱，你发出怎样的呻吟？}
\verseline{}{残酷的爱啊，你有什么不能逼迫凡人的心去做？\footnote{improbe Amor, quid non mortalia pectora cogis?复用第三卷第57行的问句；前一处谴责贪金，此处转向爱欲。}}
\verseline{}{她被迫再次流泪，再次尝试祈求，}
\verseline{}{以求告者的身份使心志向爱低下，}
\verseline{415}{免得徒然死去时还留下任何未曾尝试的事。}
\verseline{}{“安娜，你看，整个海岸都在匆忙准备；他们从}
\verseline{}{四面八方聚来；帆布已经呼唤长风，}
\verseline{}{水手们满心欢喜，已把花环放在船尾。}
\verseline{}{妹妹，既然我能够预料如此深重的痛苦，}
\verseline{420}{也就能够承受。但安娜，请为可怜的我}
\verseline{}{替我做这一件事。那个背信的人一向只}
\verseline{}{敬重你，连隐秘的心意也托付给你；}
\verseline{}{只有你能找准以温言接近那男人的门路与合宜时机。\footnote{Conte 2019 第 423 行采用 noris，并指出若读传本 noras，会暗示安娜此前已与埃涅阿斯熟识；他认为这与 mollis aditus et tempora 不合。据此，狄多是在说安娜“能找准机会，以温言软化埃涅阿斯”。}}
\verseline{}{去吧，妹妹，以求告者的身份向高傲的敌人说：}
\verseline{425}{我不曾在奥里斯与达那俄斯人一起起誓}
\verseline{}{毁灭特洛伊一族，也不曾派舰队攻打帕加马；}
\verseline{}{我没有掘动他父亲安喀塞斯的骨灰或亡魂，\footnote{奥里斯是希腊联军出征特洛伊前集结、盟誓并献祭伊菲革涅亚的地点。狄多强调自己既未参加达那俄斯人的战争誓约，也未侵犯安喀塞斯的坟墓；扰动死者的骨灰与亡魂在古代宗教观念中属于严重亵渎。}}
\verseline{}{为何他拒绝让我的话落进坚硬的耳朵？}
\verseline{}{他奔向哪里？让他给可怜的爱人这最后的赠礼：}
\verseline{430}{等到航行容易，等到长风能够托送。}
\verseline{}{我已不再祈求他背叛的旧日婚姻，}
\verseline{}{也不求他放弃美丽的拉提乌姆与王国；}
\verseline{}{我只求一点空出的时间，让狂乱稍得休息和余地，}
\verseline{}{直到命运教会战败的我怎样悲哀。}
\verseline{435}{我只求这最后的宽宥——可怜你的姐姐——}
\verseline{}{等他给我以后，我会用死亡加倍偿还。”}
\verseline{}{她这样祈求，最不幸的妹妹便把这样的哭诉}
\verseline{}{带去又带回；可任何眼泪都不能使他动摇，}
\verseline{}{任何声音都不能使他变得柔顺；}
\verseline{440}{命运横加阻挡，神明封住那男人温和的双耳。}
\verseline{}{正如阿尔卑斯的北风轮番从这一边、那一边}
\verseline{}{吹来，争着拔起一棵木质坚韧的百年橡树；}
\verseline{}{尖啸响起，高处的枝叶因树干震动}
\verseline{}{铺满大地；树本身却咬住岩石，}
\verseline{445}{树冠伸入天上空气有多高，}
\verseline{}{根须便向塔耳塔洛斯伸下多深。}
\verseline{}{英雄也这样承受两面不断的言语捶击，}
\verseline{}{在宽阔胸中感到忧苦；}
\verseline{}{心志依然不动，眼泪徒然滚落。\footnote{阿尔卑斯橡树的树冠伸向天空，根则深入塔耳塔洛斯，使埃涅阿斯的坚守具有贯通宇宙上下层的尺度。末句“眼泪徒然滚落”没有明说眼泪属于谁，通常承接埃涅阿斯；他受到触动，却仍没有改变离去的决定。}}
\verseline{450}{这时，不幸的狄多因自己的命运而惊恐，}
\verseline{}{祈求一死；她已厌倦仰望苍穹。}
\verseline{}{更催她完成既定的计划、离开光明的，是这些异兆：}
\verseline{}{她在燃香的祭坛上摆放祭品时看见——}
\verseline{}{说来也骇人——神圣的祭水变黑，}
\verseline{455}{浇下的酒化作污浊的血。}
\verseline{}{她没有把这景象告诉任何人，连安娜也没有。}
\verseline{}{宫中还有一座以大理石建成的神殿，}
\verseline{}{属于她从前的丈夫，她以非凡的敬意守护，}
\verseline{}{用雪白的羊毛与节庆的枝叶围住。}
\verseline{460}{每当幽暗的夜笼罩大地，她仿佛从那里听见}
\verseline{}{声音，听见丈夫的话在呼唤；}
\verseline{}{一只猫头鹰独自在屋顶以送葬的歌声}
\verseline{}{反复悲鸣，把长音拖成哭泣。\footnote{祭水变黑、酒变成血、亡夫神殿中的声音和屋顶猫头鹰，共同组成死亡预兆。猫头鹰的长鸣在古代占兆传统中常与凶讯、葬礼和夜间不祥相联。狄多未告诉任何人、包括妹妹的，是祭水与酒的异变（456行）。}}
\verseline{}{旧日先知的许多预言也以}
\verseline{465}{可怖的警告使她惊恐。梦中，凶狠的埃涅阿斯}
\verseline{}{亲自驱迫狂乱的她；她总觉得自己被遗下，}
\verseline{}{总觉得独自一人走在漫长的路上，没有同伴，}
\verseline{}{在荒芜的土地上寻找推罗人。}
\verseline{}{如彭透斯失去心神，望见成队的复仇女神，}
\verseline{470}{看见两个太阳、两座底比斯显现；}
\verseline{}{又如舞台上的阿伽门农之子俄瑞斯忒斯}
\verseline{}{被驱迫，逃避母亲——她手持火炬，缠着黑蛇——}
\verseline{}{而惩罚罪行的复仇女神已经坐在门前。\footnote{彭透斯在欧里庇得斯《酒神的女信徒》中因疯狂看见两个太阳、两座底比斯；俄瑞斯忒斯弑母后受复仇女神追逐。“舞台上的”一词直接让人想到悲剧演出。}}
\verseline{}{于是，痛苦征服了她，狂怒进入她心中；}
\verseline{475}{她决意死去，独自在心里选择时间和方式，}
\verseline{}{又向悲哀的妹妹开口，}
\verseline{}{用面容掩盖计划，装出怀有希望的明朗神色：}
\verseline{}{“妹妹，我找到了一个办法——为姐姐欢喜吧——}
\verseline{}{它会把他还给我，或使爱着他的我得到解脱。}
\verseline{480}{大洋的尽头、太阳沉落的近旁，}
\verseline{}{是埃塞俄比亚人最遥远的土地，至大的阿特拉斯在那里}
\verseline{}{以肩转动镶嵌着燃烧星辰的天穹。}
\verseline{}{有人从那里为我指出一位马叙利亚族的女祭司，}
\verseline{}{她守护赫斯珀里得斯的神殿，给巨龙}
\verseline{485}{喂食，看守树上的神圣枝条，}
\verseline{}{洒下湿润的蜂蜜和催眠的罂粟。\footnote{马叙利亚人是北非族群；女祭司被安置在神话中的世界西端、阿特拉斯与赫斯珀里得斯所在之地。赫斯珀里得斯守护金苹果，巨龙拉冬看守神树；蜂蜜与催眠的罂粟是这里所述祭仪的材料。狄多声称借此祭仪挽回或解除爱欲，以魔法计划掩饰自尽准备。}}
\verseline{}{她许诺用咒歌解开任何她愿意解开的心，}
\verseline{}{也把严酷的忧苦送进另一些心里；}
\verseline{}{她能使河水停住，使星辰倒转，}
\verseline{490}{还能惊动夜间的亡魂；你会看见大地}
\verseline{}{在脚下吼叫，看见白蜡树从山上走下。}
\verseline{}{亲爱的，我以众神、以你、以你甜美的生命}
\verseline{}{作证：妹妹，我并非情愿才束起身来施行魔法。}
\verseline{}{你在宫室深处秘密搭起一座露天火堆；}
\verseline{495}{那不义的人留在寝宫、挂在墙上的兵器，}
\verseline{}{他的一切遗物，还有使我毁灭的}
\verseline{}{婚床，都放在上面；我决定抹去}
\verseline{}{那个可憎男人留下的一切痕迹，女祭司为我指明办法。”}
\verseline{}{她说完便沉默；苍白同时占据面容。}
\verseline{500}{但安娜不相信姐姐正以新的祭仪}
\verseline{}{掩藏死亡，也想不到她心中有如此狂怒，}
\verseline{}{不曾担心会有比叙凯乌斯之死更沉重的事；}
\verseline{}{于是便准备所受的吩咐。}
\verseline{}{宫室深处，巨大的露天火堆}
\verseline{505}{用松木火炬和劈开的圣栎木搭成；}
\verseline{}{王后用花环围住那地方，覆以送葬的枝叶；}
\verseline{}{她把遗物和留下的剑放在上面，}
\verseline{}{又把他的像安置在卧榻上，清楚知道未来。}
\verseline{}{祭坛环立四周；女祭司披散长发，}
\verseline{510}{以雷鸣般的声音呼唤三百位神明，呼唤厄瑞玻斯与混沌，}
\verseline{}{三身的赫卡忒和处女狄安娜的三副面孔。\footnote{赫卡忒与夜、岔路、魔法和地下世界相联；她常与地上的狄安娜、天上的月神构成三重神格，因此出现“三身”和“三副面孔”。仪式使用的阿韦尔努斯泉水带有 simulati（冒充、伪称）的限定。}}
\verseline{}{她又洒下伪称取自阿韦尔努斯泉的水；}
\verseline{}{又寻来月光下用青铜镰刀割下的}
\verseline{}{饱满的药草，其中流着黑色毒乳；}
\verseline{515}{还要寻找从初生马驹额头扯下、}
\verseline{}{从母马那里抢走的爱物。\footnote{从初生马驹额头夺下的“爱物”通常与古代称作 hippomanes 相联：传说母马会立刻吞下这种肉质附着物，若被夺走便拒绝哺乳，因而成为爱情魔法的材料（见\href{https://penelope.uchicago.edu/Thayer/L/Roman/Texts/Pliny_the_Elder/8*.html}{老普林尼《博物志》8.165}）。狄多赤一足、解开衣结、持祭粉立于祭坛旁，也属于解除日常束缚的魔法仪式姿态。}}
\verseline{}{狄多自己在祭坛旁，亲手持着祭粉，伸出虔敬的双手，}
\verseline{}{一只脚解开束带，衣袍也松开腰结；}
\verseline{}{这个将死的人向众神和知晓命运的}
\verseline{520}{群星作证；若有公正而不忘旧事的神明}
\verseline{}{眷顾不平等盟约中的爱人，她便向那神祈求。\footnote{si quod…numen意为若有这样的神明：狄多祈求公正、记得旧事并眷顾不平等盟约中爱人的神。non aequo foedere可指感情不等、权力不等或盟约本身不公。诗句未指名某位神，双方关系的名分仍然悬而未决。}}
\verseline{}{那是夜晚，疲惫的身体遍布大地，}
\verseline{}{享受安宁的睡眠；森林和凶猛的海面也已平静。}
\verseline{}{群星正转过运行的中途，}
\verseline{525}{一切田野都沉默，牲畜和彩羽的鸟儿，}
\verseline{}{无论栖居在广阔清澈的湖上，还是住在遍生}
\verseline{}{荆棘的荒野，都已在寂静夜色下躺入睡眠，\footnote{众生安睡、女主人公独自不能入眠的夜景承接阿波罗尼俄斯《阿尔戈英雄纪》3.744 以下美狄亚彻夜难眠的场面。Conte 2019 正文由第 527 行直接接第 529 行；该版校勘记注明传统第 528 行在多种重要传本中缺失。}}
\verseline{}{可腓尼基女人不幸的心不能如此，她从未}
\verseline{530}{被睡眠解开，眼睛和胸膛都不接纳黑夜；}
\verseline{}{忧虑成倍增长，爱又一次升起，}
\verseline{}{狂暴起来，在愤怒的巨浪中翻涌。}
\verseline{}{她就这样继续，在心里反复思量：}
\verseline{}{“看啊，我该怎么办？再去试探从前的求婚者，}
\verseline{535}{任人嘲笑？以求告者的身份向努米底亚人求取婚姻，}
\verseline{}{向那些曾被我一再鄙弃的丈夫屈身？}
\verseline{}{那么就跟随伊利乌姆的舰队，服从特乌克罗斯人}
\verseline{}{最后的命令？难道就因为他们感激我从前的援助，}
\verseline{}{因为旧日恩惠在善于记取的人那里仍有分量？}
\verseline{540}{可即使我愿意，谁会准许，谁会在高傲的船上}
\verseline{}{接纳可憎的我？唉，迷失的人，你还不知道，}
\verseline{}{还看不清拉俄墨冬一族的背誓吗？\footnote{特洛伊王拉俄墨冬曾答应酬谢阿波罗与涅普顿修筑城墙，事后拒付报酬；又背弃对赫尔库勒斯的承诺。其“背誓”成为特洛伊王族不守信义的传统标记。狄多把祖先的罪归到埃涅阿斯及其族人身上。}}
\verseline{}{那又怎样？独自逃走，陪伴欢欣的水手，}
\verseline{}{还是由推罗人和我的全部人马簇拥}
\verseline{545}{进发，再把那些我好不容易才从西顿城拔出的人}
\verseline{}{重新赶上海，命他们把帆交给长风？}
\verseline{}{不如死吧，这是你应得的；用剑使痛苦转开。}
\verseline{}{妹妹，是你被我的眼泪征服，是你最先把狂乱的我}
\verseline{}{压进这些灾祸，抛给敌人。}
\verseline{550}{我为什么不能远离婚床，清白无咎地度过一生，}
\verseline{}{如一头野兽，不曾触及这样的忧虑！}
\verseline{}{我未能守住对叙凯乌斯的骨灰许下的信义！”}
\verseline{}{她从胸中迸发出这样沉重的悲诉。}
\verseline{}{埃涅阿斯已经决意离去，在高高的船尾}
\verseline{555}{享受睡眠，一切都已依礼备妥。}
\verseline{}{神明的形影带着先前同一副面容}
\verseline{}{在梦中向他显现，仿佛又一次警告——}
\verseline{}{一切都同墨丘利相似：声音、肤色、}
\verseline{}{金黄的头发和青春优美的肢体：}
\verseline{560}{“女神之子，在这样的危局中你还能沉入睡眠，}
\verseline{}{竟看不见往后有什么危险环伺，}
\verseline{}{失心的人，也听不见有利的西风正在吹拂？}
\verseline{}{她在胸中转动欺计和可怕之事，}
\verseline{}{决意死去，在胸中掀起一阵阵变幻的怒潮。}
\verseline{565}{趁还能一头逃走，你还不赶快离开这里？}
\verseline{}{很快你会看见船只搅乱海面，看见凶狠的}
\verseline{}{火炬闪耀，海岸因火焰沸腾，}
\verseline{}{若奥罗拉来到这片土地，撞见你仍在迟延。}
\verseline{}{来吧，行动起来，打破拖延。女人永远}
\verseline{570}{反复而易变。”说完，他融入黑夜。\footnote{varium et mutabile semper femina 是梦中神使催促埃涅阿斯立即离开时所说的话。墨丘利第二次催离的场景承接《伊利亚特》24.677—695：赫尔墨斯再次在夜间催促普里阿摩斯离开阿喀琉斯的营帐。}}
\verseline{}{埃涅阿斯被突然出现的形影惊骇，}
\verseline{}{猛然从睡梦中起身，催促同伴赶紧行动：}
\verseline{}{“醒来，男人们，坐上桨架；}
\verseline{}{立刻解开船帆。一位从高天派来的神明}
\verseline{575}{看啊，又在催我们赶快逃离，割断绞紧的缆绳。}
\verseline{}{神圣者啊，无论你是哪一位神，我们都跟随你，}
\verseline{}{欣然再次服从你的命令。}
\verseline{}{愿你在场，和善地相助，在天空带来}
\verseline{}{有利的星辰。”他说完，从鞘中拔出}
\verseline{580}{闪电般的剑，用出鞘的刀锋砍断缆索。}
\verseline{}{同样的热望同时攫住众人：他们抓起用具，奔向船只；}
\verseline{}{舰队离岸而去，海面隐没在船阵之下；}
\verseline{}{一齐划开泡沫，扫过青蓝的海面。}
\verseline{}{初晓时，奥罗拉正离开提托诺斯藏红色的床榻，}
\verseline{585}{以新生的光洒遍大地。\footnote{提托诺斯是奥罗拉的丈夫。女神为他求得永生，却未求永恒青春，于是他不断衰老； “藏红色的床榻”是史诗黎明公式的一种华丽变体，也把埃涅阿斯成功逃离的清晨同狄多即将发现空港的时刻连在一起。}}
\verseline{}{王后一从望楼看见晨光发白，}
\verseline{}{看见舰队扬着齐整的船帆前进，}
\verseline{}{觉察海岸与港湾空荡，不见一个桨手，}
\verseline{}{便用手三次、四次捶打美丽的胸膛，}
\verseline{590}{扯断金黄的头发：“朱庇特啊，”她说，“这个外来者真要走了，}
\verseline{}{嘲弄我的王国以后扬长而去？}
\verseline{}{他们还不取出兵器，从全城出发追赶？}
\verseline{}{另一些人还不从船坞扯出船只？去，}
\verseline{}{赶快拿来火把，递上兵器，奋力划桨！——}
\verseline{595}{我在说什么？我在哪里？怎样的狂乱改变了我的心？}
\verseline{}{不幸的狄多，如今这些不义之事才触到你吗？}
\verseline{}{在你交出权杖时便该如此！——看啊，那人的右手与信义，}
\verseline{}{人们说他随身带着祖先的家神，}
\verseline{}{说他曾把年迈的父亲背上双肩！}
\verseline{600}{我不能夺来他的身体，将它撕碎，抛散在}
\verseline{}{浪涛里吗？不能用剑杀死同伴，连阿斯卡尼乌斯}
\verseline{}{本人也杀死，摆在父亲的餐桌上供他吞食吗？——\footnote{狄多想象杀死阿斯卡尼乌斯并把他的肉摆给父亲。最接近的神话先例是阿特柔斯杀死堤厄斯忒斯的儿子，再设宴让父亲食用。这些报复都没有实施；她随后把暴力转向自己。}}
\verseline{}{可是战斗的胜负原本难料——难料又怎样！}
\verseline{}{一个将死的人还怕谁？我本该把火炬带进营地，}
\verseline{605}{让烈焰填满船舱，把儿子、父亲}
\verseline{}{连同整个种族消灭，再把我自己也投进去。}
\verseline{}{太阳啊，你以火光巡照世间万事，}
\verseline{}{还有你，这苦恋的媒人和知情者朱诺，}
\verseline{}{以及夜间在城市三岔路口受人号叫呼求的赫卡忒，}
\verseline{610}{惩罚罪行的复仇女神，以及看顾将死厄利莎的众神，\footnote{“厄利莎”是狄多的另一名字，古代传统通常把前者视为她原有的腓尼基名，把“狄多”解释为后来所得称号。赫卡忒在三岔路口受夜间呼求；复仇女神负责追究亲族流血和誓约破坏。狄多在这里用自己的腓尼基名“厄利莎”称呼将死的自己。}}
\verseline{}{请接受这些，让应有的神意转向我的苦难，}
\verseline{}{听取我的祈求。倘若那可憎的人必须抵达}
\verseline{}{港口，游上陆地，}
\verseline{}{倘若朱庇特的命运如此要求，这条终局已经钉牢，}
\verseline{615}{愿他受一个勇悍民族的战争和兵器折磨，}
\verseline{}{被逐出疆土，从尤卢斯的怀抱中撕开；}
\verseline{}{愿他求告援助，亲眼看见自己人的}
\verseline{}{不堪之死；即使屈身接受不平等的和平法令，}
\verseline{}{也不得享有王国和盼望的光明，}
\verseline{620}{却在死期前倒下，无人掩埋，横在沙滩中央。\footnote{诅咒先承认朱庇特的命运不可移动，再在这一前提下要求埃涅阿斯为抵达意大利付出最大代价：战争、流亡、与儿子分离、目睹同伴死亡、屈从不利和平以及早死无葬。其中若干内容预示后六卷的意大利战争； “无人掩埋”等细节则留在诗外传统或保持悬而未决。}}
\verseline{}{这就是我的祈求，是我随鲜血倾出的最后话语。}
\verseline{}{然后你们，推罗人啊，以仇恨逼迫他的后裔}
\verseline{}{和整个未来种族，把这份供物献给我的骨灰。}
\verseline{}{两族之间永无爱，也永无盟约。}
\verseline{625}{愿有一位复仇者从我的尸骨中兴起，}
\verseline{}{用火炬与剑追逐达耳达诺斯人的殖民者，}
\verseline{}{无论现在、将来，无论何时力量允许。}
\verseline{}{我诅咒海岸敌对海岸，浪涛敌对浪涛，}
\verseline{}{兵器敌对兵器；让他们自己和他们的后代永远争战。”\footnote{“从我的尸骨中升起吧，一个复仇者”没有点名具体人物；对奥古斯都时代的罗马读者而言，它最强烈地唤起汉尼拔和布匿战争。狄多诅咒“海岸敌对海岸、浪涛敌对浪涛、兵器敌对兵器”，并祈求后代永远交战；诗由此把她的诅咒同后来罗马与迦太基的战争联系起来。}}
\verseline{630}{她说完，心思仍转向各个方向，}
\verseline{}{寻找尽快割断这可憎光明的办法。}
\verseline{}{接着，她简短吩咐叙凯乌斯的乳母巴耳刻；}
\verseline{}{她自己的乳母已化成黑灰，留在古老祖国：}
\verseline{}{“亲爱的乳母，把我妹妹安娜带到这里；}
\verseline{635}{叫她赶快用河水洒净身体，}
\verseline{}{带来祭牲和已经指明的赎罪祭物；}
\verseline{}{让她这样前来；你自己也用神圣的祭带覆住额头。}
\verseline{}{我已依礼开始为冥府的朱庇特准备祭仪，}
\verseline{}{如今决意将它完成，使忧苦终止，}
\verseline{640}{让火焰吞没那个达耳达诺斯人的火堆。”\footnote{“冥府的朱庇特”是地下世界统治者狄斯的称号。狄多以烧毁埃涅阿斯遗物为名支开安娜和乳母，堆起的火台最终成为她自己的葬火。}}
\verseline{}{她说完，老妇便以老年人的急切加快脚步。}
\verseline{}{狄多却因这骇人的举动战栗、狂烈，}
\verseline{}{转动血红的目光，颤抖的双颊}
\verseline{}{泛起斑痕，又因将临的死亡而苍白。}
\verseline{645}{她冲破宫中内室的门槛，}
\verseline{}{狂乱地登上高高的火堆，拔出那把达耳达诺斯人的}
\verseline{}{剑——这件礼物并非为这样的用途求来。\footnote{狄多用埃涅阿斯留下的剑自尽。原文说这件礼物“原不曾为这样的用途求来”，宾客赠礼或爱人遗物遂反转为死亡工具；火堆上的衣袍、卧榻、画像与剑也把私人记忆组织成一场替身葬礼。}}
\verseline{}{她在那里看见伊利乌姆的衣袍和熟悉的卧榻，}
\verseline{}{泪水与思绪使她停留片刻，}
\verseline{650}{她便伏在榻上，说出最后的话：}
\verseline{}{“遗物啊，当命运和神明还容许时，你们曾甜美，}
\verseline{}{请接纳这个灵魂，使我从这些忧苦中解脱。}
\verseline{}{我活过，也走完了命运给我的路程，}
\verseline{}{如今我伟大的幽影将走入地下。}
\verseline{655}{我建立一座显赫的城，看见自己的墙垣；}
\verseline{}{我为丈夫复仇，也惩罚了敌对的兄弟；}
\verseline{}{我本来幸福，唉，幸福至极——只要达耳达诺斯人的}
\verseline{}{船从未触及我们的海岸！”\footnote{狄多临死前以近似罗马政治人物墓志铭的方式回顾自己：活过既定路程、建立名城、看见城墙、为丈夫复仇、惩罚敌对的兄弟。她对皮格马利翁的“惩罚”指携带财富逃离，并挫败他的谋财与控制。卡图卢斯第 64 首的阿里阿德涅也在绝望中祈愿雅典船从未抵达克里特；马克罗比乌斯（6.1.42）明确把两处诗句并列。}}
\verseline{}{她说完，把唇压在卧榻上：“我将死而未得复仇，}
\verseline{660}{可还是让我死吧。就这样，就这样走进幽影，我心甘情愿。}
\verseline{}{让那残忍的达耳达诺斯人在远海把这火光饮进眼中，}
\verseline{}{带着我死亡的凶兆上路。”}
\verseline{}{她说完。侍女们正听着这些话，忽见她}
\verseline{}{倒在剑上；剑刃因鲜血冒着泡沫，}
\verseline{665}{双手也被溅满。呼声冲进高耸的}
\verseline{}{厅堂；流言在震动的全城发起狂舞。}
\verseline{}{屋宇因悲号、呻吟和女人的哀叫}
\verseline{}{震响；天空回荡着捶胸的巨声，}
\verseline{}{仿佛敌军已经闯入，整个迦太基}
\verseline{670}{或古老的推罗正在陷落，狂暴的火焰}
\verseline{}{滚过人类与神明的屋顶。}
\verseline{}{妹妹听见，惊得失魂，急促奔来，}
\verseline{}{妹妹以指甲毁伤面容，以拳头捶打胸膛，}
\verseline{}{冲过人群，呼喊垂死者的名字：}
\verseline{675}{“姐姐，原来是这样？你竟用欺骗对付我？}
\verseline{}{这座火堆、这些火焰和祭坛，竟为我准备这样的结局？}
\verseline{}{我被遗弃，先该悲叹什么？你临死时}
\verseline{}{竟轻视妹妹，不让她相伴？你本该叫我走向同一命运：}
\verseline{}{同样的痛苦，同一个时刻，会用同一把剑带走我们。}
\verseline{680}{我还亲手搭起这座火堆，亲口呼唤}
\verseline{}{祖先的神明，到头来，你这样躺着，我却不在身旁——何等残酷！\footnote{crudelis（“残酷的”）既可指狄多，也可指安娜自己：前者是怨姐姐弃她而死，后者是自责曾帮助姐姐筑起葬火。古代注家已有这两种解释。}}
\verseline{}{姐姐，你毁了自己，也毁了我，毁了人民、}
\verseline{}{西顿的长老和你的城。容我用清水洗净}
\verseline{}{伤口；倘若最后一息还在她身上游荡，}
\verseline{685}{让我用自己的嘴接住。”她这样说着，已经登上高高的台阶，\footnote{安娜希望用自己的嘴接住狄多唇边游荡的最后一口气。古代文学常把生命想象成随气息或灵魂从口中逸出；亲属贴近死者嘴唇，既是最后亲吻，也可被理解为接取临终气息的动作。Conte 2019 校勘记将此处的 date 解作 sinite（“容我”）：安娜请求获准用水洗伤，并接住最后一息。}}
\verseline{}{把半死的姐姐抱在胸前，用怀抱温暖，}
\verseline{}{一边呻吟，一边用衣袍擦去黑色的血。}
\verseline{}{狄多试图抬起沉重的眼睛，又一次}
\verseline{}{失去力量；钉在胸下的伤口嘶嘶作响。}
\verseline{690}{她三次用手肘撑起身来；}
\verseline{}{三次又倒回卧榻，游移的眼睛在高深的}
\verseline{}{天空寻找光明，一找到便发出呻吟。}
\verseline{}{全能的朱诺怜悯她漫长的痛苦}
\verseline{}{和艰难的死亡，便从奥林匹斯派下伊里斯，}
\verseline{695}{解开挣扎的灵魂和被缚住的肢体。}
\verseline{}{她并非在命定之时死去，也不是罪有应得地死去；}
\verseline{}{可怜的她未至死期便被骤起的狂怒点燃，}
\verseline{}{普洛塞庇娜还未从她头顶取走金黄的}
\verseline{}{发丝，还未把她本人判给斯提克斯的奥尔库斯。}
\verseline{700}{于是伊里斯用沾露的藏红羽翼穿过天空，}
\verseline{}{迎着太阳拖出千种变幻的颜色，}
\verseline{}{飞落下来，停在她头顶：“我奉命把这缕}
\verseline{}{圣物带给狄斯，让你从这具身体解脱。”}
\verseline{}{她说完，右手剪下发丝；所有温热顿时}
\verseline{705}{消散，生命退进风中。\footnote{狄多三次试图起身、三次倒下；三重动作的节奏也见于第二卷埃涅阿斯三次试图拥抱克瑞乌萨的场面。她死于命定时刻之前，这死亡也并非她所应得；普洛塞庇娜尚未剪取应献给地下世界的发丝，所以灵魂迟迟不能脱离身体。朱诺派伊里斯代行剪发；伊里斯沾露的藏红羽翼在阳光中显出虹霓，明亮的色彩落在狄多身体温热消散的最后一刻。}}
